\chapter{The playground}
\label{ch:playground}

Before Scheme meets files, models and guards, let's play. Every program in
this chapter runs in \code{nodwyn repl} or with \code{nodwyn eval FILE}, and
each one shows off a technique you will use later on real work.

\section{FizzBuzz, the traditional warm-up}

\begin{scm}
(define (fizzbuzz n)
  (cond ((= 0 (remainder n 15)) "FizzBuzz")
        ((= 0 (remainder n 3)) "Fizz")
        ((= 0 (remainder n 5)) "Buzz")
        (else (number->string n))))

(display (string-join (map fizzbuzz (cdr (iota 16))) " "))
\end{scm}

\begin{txt}
1 2 Fizz 4 Buzz Fizz 7 8 Fizz Buzz 11 Fizz 13 14 FizzBuzz
\end{txt}

Technique: \code{map} over \code{iota} instead of a loop, and
\code{string-join} to glue the results.

\section{The Towers of Hanoi}

Move a stack of discs from peg \code{a} to peg \code{c}, one at a time, never
putting a bigger disc on a smaller one. The recursive answer is famously
short: move all but the bottom disc out of the way, move the bottom one,
move the rest back on top.

\begin{scm}
(define (hanoi n from to via)
  (if (= n 0)
      '()
      (append (hanoi (- n 1) from via to)
              (list (list from to))
              (hanoi (- n 1) via to from))))

(hanoi 3 'a 'c 'b)
;=> ((a c) (a b) (c b) (a c) (b a) (b c) (a c))
(length (hanoi 10 'a 'c 'b))      ;=> 1023
\end{scm}

Technique: the answer is \emph{data}, a list of moves, rather than
printing. Data can be counted, checked and turned into text later. That is
the habit that makes strategies testable.

\section{The sieve of Eratosthenes}

A vector of flags, crossed off in multiples:

\begin{scm}
(define (primes-up-to n)
  (let ((sieve (list->vector (map (lambda (i) #t) (iota (+ n 1))))))
    (let mark ((i 2))
      (when (<= (* i i) n)
        (when (vector-ref sieve i)
          (do ((j (* i i) (+ j i))) ((> j n))
            (vector-set! sieve j #f)))
        (mark (+ i 1))))
    (filter (lambda (i) (and (>= i 2) (vector-ref sieve i)))
            (iota (+ n 1)))))

(primes-up-to 50)
;=> (2 3 5 7 11 13 17 19 23 29 31 37 41 43 47)
\end{scm}

Technique: a vector for random access, \code{do} for a counting loop, a
named \code{let} for the outer loop, and a list for the answer.

\section{Counting words}

Which words does a text use most? Build an association list of counts with
\code{fold}, then sort it:

\begin{scm}
(define (count-words text)
  (fold (lambda (counts word)
          (if (assoc word counts)
              (map (lambda (p)
                     (if (equal? (car p) word) (cons word (+ 1 (cdr p))) p))
                   counts)
              (cons (cons word 1) counts)))
        '()
        (filter (lambda (w) (not (string-empty? w)))
                (string-split (string-downcase text) " "))))

(define (take l n)
  (if (or (null? l) (= n 0)) '() (cons (car l) (take (cdr l) (- n 1)))))

(define (top-words text n)
  (take (sort (count-words text) (lambda (a b) (> (cdr a) (cdr b)))) n))

(top-words "the basil and the tomato and the bee" 2)
;=> (("the" . 3) ("and" . 2))
\end{scm}

Technique: an association list as a small dictionary, rebuilt rather than
changed. \nodwyn's optimizer (Chapter~\ref{ch:habits}) digests a session's
log with exactly this kind of code.

\section{Roman numerals}

\begin{scm}
(define roman-table
  '((1000 . "M") (900 . "CM") (500 . "D") (400 . "CD") (100 . "C")
    (90 . "XC") (50 . "L") (40 . "XL") (10 . "X") (9 . "IX")
    (5 . "V") (4 . "IV") (1 . "I")))

(define (roman n)
  (let loop ((n n) (table roman-table) (out ""))
    (cond ((= n 0) out)
          ((>= n (caar table))
           (loop (- n (caar table)) table (string-append out (cdar table))))
          (else (loop n (cdr table) out)))))

(map roman '(4 9 14 1990 2026))
;=> ("IV" "IX" "XIV" "MCMXC" "MMXXVI")
\end{scm}

Technique: a table as data and a tail-recursive loop that carries three
things at once. \code{caar} is ``the car of the car'': the number of the
first pair in the table.

\section{A calculator: Scheme reading Scheme}

Because code is lists, a program can read a quoted expression and
evaluate it. Here is a calculator for arithmetic with variables:

\begin{scm}
(define (calc expr env)
  (cond ((number? expr) expr)
        ((symbol? expr)
         (let ((b (assq expr env)))
           (if b (cdr b) (error "calc: unknown variable" expr))))
        ((pair? expr)
         (let ((op (car expr))
               (args (map (lambda (e) (calc e env)) (cdr expr))))
           (case op
             ((+) (apply + args))
             ((-) (apply - args))
             ((*) (apply * args))
             ((/) (apply / args))
             ((max) (apply max args))
             (else (error "calc: unknown operator" op)))))
        (else (error "calc: cannot evaluate" expr))))

(calc '(+ 1 (* 2 3)) '())                           ;=> 7
(calc '(* price (+ 1 tax)) '((price . 20) (tax . 0.25)))  ;=> 25.0
(guard (e (#t (assoc-ref e 'message)))
  (calc '(sqrt 2) '()))                             ;=> "calc: unknown operator"
\end{scm}

Look at what \code{calc} does: it decides \emph{what each operator is
allowed to do}, and refuses the rest. That is a guard in miniature. Keep it
in mind for Chapter~\ref{ch:primitives}.

\section{Memory: Fibonacci without the wait}

The naive Fibonacci recomputes the same values exponentially often. Keep
each answer in a vector the first time:

\begin{scm}
(define (make-memo size) (list->vector (map (lambda (i) #f) (iota size))))

(define (fib n)
  (let ((memo (make-memo (+ n 1))))
    (let f ((k n))
      (cond ((< k 2) k)
            ((vector-ref memo k) => (lambda (v) v))
            (else (let ((v (+ (f (- k 1)) (f (- k 2)))))
                    (vector-set! memo k v)
                    v))))))

(fib 90)                        ;=> 2880067194370816120
\end{scm}

Technique: \code{cond}'s \code{=>} passes the test's value on, so ``if we
have it, return it'' is one clause.

\section{Rolling dice without a random number}

There is no \code{random} in the language, and that is on purpose:
strategies should be reproducible from the log. But a closure and some
arithmetic make a perfectly good repeatable die:

\begin{scm}
(define (make-dice seed)
  (lambda (sides)
    (set! seed (modulo (+ (* seed 1103515245) 12345) 2147483648))
    (+ 1 (remainder seed sides))))

(define d6 (make-dice 42))
(map (lambda (i) (d6 6)) (iota 8))     ;=> (4 3 6 5 2 1 6 1)
\end{scm}

The same seed gives the same rolls every time: a linear congruential
generator, the kind early computers used.

\section{Drawing a tree}

\begin{scm}
(define garden
  '("garden" ("herbs" ("basil") ("mint")) ("veg" ("tomato") ("bean"))))

(define (show-tree tree depth)
  (display (string-append (substring "            " 0 (* 2 depth)) (car tree)))
  (newline)
  (for-each (lambda (child) (show-tree child (+ depth 1))) (cdr tree)))

(show-tree garden 0)
\end{scm}

\begin{txt}
garden
  herbs
    basil
    mint
  veg
    tomato
    bean
\end{txt}

\section{A Markdown table}

Models read Markdown well, and a tool that answers with a table is easier
for a model to use than raw data:

\begin{scm}
(define (md-table header rows)
  (define (row cells) (string-append "| " (string-join cells " | ") " |"))
  (string-join
   (append (list (row header) (row (map (lambda (h) "---") header)))
           (map (lambda (r) (row (map (lambda (x) (format "~a" x)) r))) rows))
   "\n"))

(display (md-table '("plant" "water") '(("basil" "daily") ("tomato" "2 days"))))
\end{scm}

\begin{txt}
| plant | water |
| --- | --- |
| basil | daily |
| tomato | 2 days |
\end{txt}

You will register a tool built on this in Chapter~\ref{ch:tools}.

\begin{puzzle}[ 10: extend the calculator]
Teach \code{calc} two new forms: \code{(if test then else)} where
\code{test} is true when it is not zero, and \code{(let ((name expr))
body)} with one binding. \code{(calc '(let ((x 4)) (if (- x 4) 1 2)) '())}
should be \code{2}.\ans{10}
\end{puzzle}

\begin{puzzle}[ 11: Hanoi, counted]
Without building the list of moves, write \code{(hanoi-count n)} that
returns how many moves \code{n} discs need, in tail position. Check it
against \code{(length (hanoi 10 'a 'c 'b))}.\ans{11}
\end{puzzle}

\begin{puzzle}[ 12: Roman numerals back]
Write \code{(unroman s)} so that \code{(unroman "MMXXVI")} is
\code{2026}. Hint: walk the same table, eating matching prefixes with
\code{string-prefix?} and \code{substring}.\ans{12}
\end{puzzle}
